\documentclass[10pt,twocoloumn]{article}

\usepackage{amsmath}
\usepackage{amsthm}
\usepackage{kida}
\usepackage{physics}
\usepackage{float}
\usepackage{svg}
\usepackage{graphicx}
\newcommand{\addsvg}[4][1.0]{
    \begin{figure}[H]
        \centering
        \includesvg[width=#1\textwidth]{#2}
        \caption{#4}
        \label{fig:#3}
    \end{figure}
}

\title{Advanced Statistical Mechanics Assignment 1}
\author{Amir Hossein Ebrahimnezhad - 404461001}
\date{due: \today}

\begin{document}
\maketitle
\section*{Problem 1}
\subsection*{Problem Statement}
In an isothermal Joule free expansion at temperature $T$, from the initial state $(P_i,V_i)$ to the final state $(P_f,V_f = 2V_i)$, with robust reasoning, calculate $\Delta U$, $\Delta H$ , $\Delta F$ and $\Delta G$. Discuss the similarities and differences.

\subsection*{Solution}

A Joule free expansion is a thermodynamic process in which a gas expands into a vacuum without doing any work and without exchanging heat with its surroundings. This process is also known as free expansion or unrestrained expansion. The key characteristics are that the expansion occurs adiabatically (no heat transfer, $Q = 0$) and against zero external pressure (no work done, $W = 0$). Despite being isothermal in the final equilibrium state for an ideal gas, the process itself is irreversible and occurs spontaneously.

\subsubsection*{Calculation of $\Delta U$ (Change in Internal Energy)}

By the first law of thermodynamics, we have:
\begin{equation}\Delta U = Q - W\end{equation}

Since the process is a free expansion into vacuum, no work is done by or on the gas ($W = 0$), and no heat is exchanged with the surroundings ($Q = 0$). Therefore:
\begin{equation}\Delta U = 0 - 0 = 0\end{equation}

For an ideal gas, the internal energy depends only on temperature. Since the process is isothermal ($T = \text{constant}$), this result is consistent with the ideal gas model. Thus, $\boxed{\Delta U = 0}$.

\subsubsection*{Calculation of $\Delta H$ (Change in Enthalpy)}

The enthalpy is defined as:
\begin{equation}H = U + PV\end{equation}

Taking the differential, we obtain:
\begin{equation}dH = dU + d(PV) = dU + PdV + VdP\end{equation}

For a finite change:
\begin{equation}\Delta H = \Delta U + \Delta(PV)\end{equation}

For an ideal gas, $PV = nRT$, so:
\begin{equation}\Delta H = \Delta U + \Delta(nRT) = \Delta U + nR\Delta T\end{equation}

Since the temperature remains constant during an isothermal process ($\Delta T = 0$), we have:
\begin{equation}\Delta H = \Delta U + 0 = 0\end{equation}

Therefore,

\begin{equation}\boxed{\Delta H = 0}.\end{equation}

\subsubsection*{Calculation of $\Delta F$ (Change in Helmholtz Free Energy)}

The Helmholtz free energy is defined as:
\begin{equation}F = U - TS\end{equation}

For an isothermal process at constant temperature $T$, the change in Helmholtz free energy is:
\begin{equation}\Delta F = \Delta U - T\Delta S\end{equation}

Since $\Delta U = 0$, we need to calculate the entropy change $\Delta S$. For an ideal gas undergoing an isothermal expansion, the entropy change is given by:
\begin{equation}\Delta S = nR\ln\left(\frac{V_f}{V_i}\right)\end{equation}

Given that $V_f = 2V_i$, we have:
\begin{equation}\Delta S = nR\ln\left(\frac{2V_i}{V_i}\right) = nR\ln(2)\end{equation}

The entropy increases because the gas has expanded into a larger volume, increasing the number of accessible microstates. Substituting into the expression for $\Delta F$:
\begin{equation}\Delta F = 0 - T \cdot nR\ln(2) = -nRT\ln(2)\end{equation}

Therefore,
\begin{equation}\boxed{\Delta F = -nRT\ln(2)}.\end{equation}

\subsubsection*{Calculation of $\Delta G$ (Change in Gibbs Free Energy)}

The Gibbs free energy is defined as:
\begin{equation}G = H - TS\end{equation}

For an isothermal process, the change in Gibbs free energy is:
\begin{equation}\Delta G = \Delta H - T\Delta S\end{equation}

Since $\Delta H = 0$ and $\Delta S = nR\ln(2)$, we have:
\begin{equation}\Delta G = 0 - T \cdot nR\ln(2) = -nRT\ln(2)\end{equation}

Therefore, $\boxed{\Delta G = -nRT\ln(2)}$.

\subsubsection*{Discussion of Similarities and Differences}

The results reveal important thermodynamic insights about the Joule free expansion process. Both the internal energy and enthalpy remain unchanged ($\Delta U = \Delta H = 0$) because the process is isothermal for an ideal gas, and both quantities depend only on temperature for ideal gases. This reflects the fact that no energy is transferred to or from the system during the expansion.

In contrast, both free energy functions decrease by the same amount ($\Delta F = \Delta G = -nRT\ln(2) < 0$). This negative change indicates that the process is spontaneous and irreversible, consistent with the second law of thermodynamics. The entropy of the system increases ($\Delta S = nR\ln(2) > 0$) as the gas occupies a larger volume, which drives the decrease in free energies.

The equality $\Delta F = \Delta G$ occurs specifically because both $\Delta U = 0$ and $\Delta H = 0$. From the definitions, we can write:
\begin{align}
    \Delta F &= \Delta U - T\Delta S = -T\Delta S \\
    \Delta G &= \Delta H - T\Delta S = -T\Delta S
\end{align}

Since $\Delta U = \Delta H$ for this process, it follows that $\Delta F = \Delta G$. This equality is particular to isothermal processes involving ideal gases where the internal energy and enthalpy changes vanish. The negative values of the free energies quantify the maximum work that could have been extracted if the expansion had been carried out reversibly, which in this irreversible free expansion was entirely dissipated as lost opportunity.
\section*{Problem 2}
\subsection*{Problem Statement}
\textit{Filament:} for an elastic filament it is found that, at a finite range in temperature, a displacement $x$ requires a force:
\begin{equation}
    J = ax - bT + cT x
\end{equation}
where $a,b$ and $c$ are constants. Furthermore, its heat capacity at constant displacement is proportional to temperature, that is:
\begin{equation}
    C_x = A(x)T
\end{equation}
\begin{itemize}
    \item Use an appropriate Maxwell relation to calculate $\partial S / \partial x |_T$.
    \item Show that $A$ has to in fact be independent of $x$, that is $dA/dx = 0$.
    \item Give the expression for $S(T,x)$ assuming $S(0,0)=S_0$.
    \item Calculate the heat capacity at constant tension, that is $C_J = T\partial S / \partial T|_J$ as a function of $T$ and $J$
\end{itemize}

\subsection*{Solution}

This problem concerns the thermodynamics of an elastic filament, which is analogous to a one-dimensional elastic system where mechanical work is done through extension or compression. The natural variables for this system are temperature $T$ and displacement $x$, with the conjugate force being the tension $J$. We will use thermodynamic relations and Maxwell relations to extract information about the entropy and heat capacity of the system.

\subsubsection*{Part 1: Maxwell Relation for $\partial S / \partial x |_T$}

For a system where work is done through displacement, the fundamental thermodynamic relation (first law) can be written as:
\begin{equation}
    dU = TdS + Jdx
\end{equation}

This suggests that the natural thermodynamic potential to use is the Helmholtz free energy $F(T,x)$, defined as:
\begin{equation}
    F = U - TS
\end{equation}

Taking the total differential:
\begin{equation}dF = dU - TdS - SdT = TdS + Jdx - TdS - SdT = -SdT + Jdx\end{equation}

From this expression, we can identify the natural variables of $F$ as $T$ and $x$, with:
\begin{equation}
    \left(\frac{\partial F}{\partial T}\right)_x = -S \quad \text{and} \quad \left(\frac{\partial F}{\partial x}\right)_T = J
\end{equation}

Since $F$ is a state function, its mixed partial derivatives must be equal (Schwarz's theorem):
\begin{equation}
    \frac{\partial^2 F}{\partial x \partial T} = \frac{\partial^2 F}{\partial T \partial x}
\end{equation}

This yields the Maxwell relation:
\begin{equation}
    \left(\frac{\partial J}{\partial T}\right)_x = -\left(\frac{\partial S}{\partial x}\right)_T
\end{equation}

Now we use the given force-displacement-temperature relation $J = ax - bT + cTx$. Taking the partial derivative with respect to $T$ at constant $x$:
\begin{equation}
    \left(\frac{\partial J}{\partial T}\right)_x = -b + cx
\end{equation}

Therefore, from the Maxwell relation:
\begin{equation}
    \boxed{\left(\frac{\partial S}{\partial x}\right)_T = b - cx}
\end{equation}

\subsubsection*{Part 2: Showing that $A$ is Independent of $x$}

The heat capacity at constant displacement is defined as:
\begin{equation}
    C_x = T\left(\frac{\partial S}{\partial T}\right)_x = A(x)T
\end{equation}

This gives us:
\begin{equation}
    \left(\frac{\partial S}{\partial T}\right)_x = A(x)
\end{equation}

Since entropy is a state function, its mixed partial derivatives must be equal:
\begin{equation}
\frac{\partial^2 S}{\partial x \partial T} = \frac{\partial^2 S}{\partial T \partial x}
\end{equation}

Taking the partial derivative of $\left(\frac{\partial S}{\partial x}\right)_T = b - cx$ with respect to $T$ at constant $x$:
\begin{equation}
    \frac{\partial}{\partial T}\left(\frac{\partial S}{\partial x}\right)_T \bigg|_x = \frac{\partial}{\partial T}(b - cx)\bigg|_x = 0
\end{equation}

Taking the partial derivative of $\left(\frac{\partial S}{\partial T}\right)_x = A(x)$ with respect to $x$ at constant $T$:

\begin{equation}
    \frac{\partial}{\partial x}\left(\frac{\partial S}{\partial T}\right)_x \bigg|_T = \frac{dA(x)}{dx}
\end{equation}
By the equality of mixed partials:
\begin{equation}
\frac{dA(x)}{dx} = 0
\end{equation}
This demonstrates that $A(x)$ must be a constant, independent of $x$. We can therefore write $A(x) = A$, where $A$ is a constant. Thus, $\boxed{\frac{dA}{dx} = 0}$.

\subsubsection*{Part 3: Expression for $S(T,x)$}

Now we have the two partial derivatives of entropy:
\begin{equation}
    \left(\frac{\partial S}{\partial T}\right)_x = A \quad \text{and} \quad \left(\frac{\partial S}{\partial x}\right)_T = b - cx
\end{equation}
To find $S(T,x)$, we integrate these expressions. Starting with the first:
\begin{equation}
    S(T,x) = \int A , dT = AT + f(x)
\end{equation}

where $f(x)$ is an arbitrary function of $x$ alone. Now we use the second partial derivative condition:
\begin{equation}
\left(\frac{\partial S}{\partial x}\right)_T = \frac{\partial}{\partial x}[AT + f(x)] = f'(x) = b - cx
\end{equation}
Integrating with respect to $x$:
\begin{equation}
    f(x) = \int (b - cx) \, dx = bx - \frac{c}{2}x^2 + \text{constant}
    \end{equation}

Therefore, the general form of the entropy is:
\begin{equation}
    S(T,x) = AT + bx - \frac{c}{2}x^2 + S_{\text{const}}
    \end{equation}
Using the boundary condition $S(0,0) = S_0$:
\begin{equation}
    S_0 = A(0) + b(0) - \frac{c}{2}(0)^2 + S_{\text{const}} = S_{\text{const}}
\end{equation}
Thus:
\begin{equation}
    \boxed{S(T,x) = S_0 + AT + bx - \frac{c}{2}x^2}
    \end{equation}
\subsubsection*{Part 4: Heat Capacity at Constant Tension $C_J$}

The heat capacity at constant tension is defined as:
\begin{equation}
    C_J = T\left(\frac{\partial S}{\partial T}\right)_J
    \end{equation}
To evaluate this, we need to express the derivative at constant $J$ rather than constant $x$. We use the general relation for changing variables:
\begin{equation}
    \left(\frac{\partial S}{\partial T}\right)_J = \left(\frac{\partial S}{\partial T}\right)_x + \left(\frac{\partial S}{\partial x}\right)_T \left(\frac{\partial x}{\partial T}\right)_J
    \end{equation}

We already know:
\begin{equation}\left(\frac{\partial S}{\partial T}\right)_x = A \quad \text{and} \quad \left(\frac{\partial S}{\partial x}\right)_T = b - cx
    \end{equation}

To find $\left(\frac{\partial x}{\partial T}\right)_J$, we use the constraint that $J$ is constant. From $J = ax - bT + cTx$, taking the total differential and setting $dJ = 0$:
\begin{align} 0 &= a\,dx - b\,dT + c(T\,dx + x\,dT) \\
&= (a + cT)dx + (cx - b)dT\end{align}

Therefore:
\begin{equation}
    \left(\frac{\partial x}{\partial T}\right)_J = -\frac{cx - b}{a + cT} = \frac{b - cx}{a + cT}
    \end{equation}

Substituting back:
\begin{equation}
    \left(\frac{\partial S}{\partial T}\right)_J = A + (b - cx) \cdot \frac{b - cx}{a + cT} = A + \frac{(b - cx)^2}{a + cT}
    \end{equation}
Therefore, the heat capacity at constant tension is:
\begin{equation}
    C_J = T\left(\frac{\partial S}{\partial T}\right)_J = AT + \frac{T(b - cx)^2}{a + cT}
    \end{equation}
To express this in terms of $T$ and $J$ rather than $x$, we solve for $x$ from the equation $J = ax - bT + cTx = x(a + cT) - bT$:
\begin{equation}
    x = \frac{J + bT}{a + cT}
    \end{equation}
Then:
\begin{equation}
    b - cx = b - c\cdot\frac{J + bT}{a + cT} = \frac{b(a + cT) - c(J + bT)}{a + cT} = \frac{ab - cJ}{a + cT}
    \end{equation}
Substituting:
\begin{equation}\boxed{C_J(T,J) = AT + \frac{T(ab - cJ)^2}{(a + cT)^3}}\end{equation}

This expression shows how the heat capacity at constant tension depends on both the temperature and the applied force, reflecting the coupling between thermal and mechanical properties of the elastic filament.
\section*{Problem 3}
\subsection*{Problem Statement}
\textit{Superconducting transition:} many metals become superconductors at low temperatures $T$, and magnetic fields $B$. The heat capacities of the two phases at zero magnetic field are approximately given by:
\begin{equation}
    \begin{cases}
        C_s(T) = V\alpha T^3 \ \ \text{in the superconducting phase}\\
        C_n(T) = V\left[\beta T^3 + \gamma T \right]\ \ \text{in the normal phase}\\
    \end{cases}
\end{equation}
where $V$ is the volume, and $\{\alpha,\beta,\gamma\}$ are constants. (There is no appreciable change in volume at this transition, and mechanical work can be ignored throughout this problem.)
\begin{itemize}
    \item Calculate the entropies $S_s(T)$ and $S_n(T)$ of the two phases at zero field, using the third law of thermodynamics.
    \item Experiments indicate that there is no latent heat ($L=0$) for the transition between the normal and superconducting phases at zero field. Use this information to obtain the transition temperature $T_c$ as a function of $\alpha, \beta$ and $\gamma$.
    \item At zero temperature, the electrons in the superconductor form bound Cooper pairs. As a result, the internal energy of superconductor is reduced by an amount $V\Delta$, that is $E_n(T=0)=E_0$ and $E_s(T=0)=E_0 -V \Delta$ for the metal and superconductor, respectively. Calculate the internal energies of both phases at finite temperatures.
    \item By comparing the Gibbs free energies (or chemical potentials) in the two phases, obtain an expression for the energy gap $\Delta$ in terms of $\alpha, \beta$ and $\gamma$.
    \item In the presence of a magnetic field $B$, inclusion of magnetic work results in $dE = TdS + BdM + \mu dN$, where $M$ is the magnetization. The superconducting phase is a perfect diamagnet, expelling the magnetic field from its interior, such that $M_s = -VB/(4\pi)$ in appropriate units. The normal metal can be regarded as approximately non-magnetic, with $M_n=0$. Use this information, in conjunction with previous results, to show that the superconducting phase becomes normal for magnetic fields larger than:
    \begin{equation}
        B_c(T) = B_0\left(1- \frac{T^2}{T_c^2}\right),
    \end{equation}
    giving an expression for $B_0$.
\end{itemize}

\subsection*{Solution}

This problem explores the thermodynamics of the superconducting phase transition, a phenomenon where certain metals lose all electrical resistance below a critical temperature. The transition between the normal and superconducting phases involves subtle changes in entropy, internal energy, and free energy that can be analyzed using fundamental thermodynamic principles.

\subsubsection*{Part 1: Entropies of the Two Phases}

The heat capacity at constant volume (or constant pressure, since volume changes are negligible) is related to entropy by:
\begin{equation}C = T\left(\frac{\partial S}{\partial T}\right) \quad \Rightarrow \quad dS = \frac{C}{T}dT\end{equation}

For the superconducting phase, we have:
\begin{equation}dS_s = \frac{C_s(T)}{T}dT = \frac{V\alpha T^3}{T}dT = V\alpha T^2 dT\end{equation}

Integrating from $T=0$ to $T$:
\begin{equation}S_s(T) = S_s(0) + \int_0^T V\alpha T'^2 dT' = S_s(0) + V\alpha \left[\frac{T'^3}{3}\right]_0^T = S_s(0) + \frac{V\alpha T^3}{3}\end{equation}

By the third law of thermodynamics, the entropy of a perfect crystalline substance approaches zero as the temperature approaches absolute zero: $S_s(0) = 0$. Therefore:
\begin{equation}\boxed{S_s(T) = \frac{V\alpha T^3}{3}}\end{equation}

For the normal phase:
\begin{equation}dS_n = \frac{C_n(T)}{T}dT = \frac{V(\beta T^3 + \gamma T)}{T}dT = V(\beta T^2 + \gamma)dT\end{equation}

Integrating from $T=0$ to $T$ with $S_n(0) = 0$:
\begin{equation}S_n(T) = \int_0^T V(\beta T'^2 + \gamma)dT' = V\left[\frac{\beta T'^3}{3} + \gamma T'\right]_0^T\end{equation}

\begin{equation}\boxed{S_n(T) = V\left(\frac{\beta T^3}{3} + \gamma T\right)}\end{equation}

\subsubsection*{Part 2: Transition Temperature $T_c$}

The latent heat of a phase transition is defined as:
\begin{equation}L = T\Delta S = T(S_n - S_s)\end{equation}

At the transition temperature $T_c$, we are told that the latent heat is zero: $L(T_c) = 0$. This means:
\begin{equation}S_n(T_c) = S_s(T_c)\end{equation}

Substituting our expressions for the entropies:
\begin{equation}V\left(\frac{\beta T_c^3}{3} + \gamma T_c\right) = \frac{V\alpha T_c^3}{3}\end{equation}

Dividing through by $V$ and multiplying by 3:
\begin{equation}\beta T_c^3 + 3\gamma T_c = \alpha T_c^3\end{equation}

Rearranging:
\begin{equation}(\alpha - \beta)T_c^3 = 3\gamma T_c\end{equation}

Since $T_c \neq 0$, we can divide by $T_c$:
\begin{equation}(\alpha - \beta)T_c^2 = 3\gamma\end{equation}

Therefore:
\begin{equation}\boxed{T_c = \sqrt{\frac{3\gamma}{\alpha - \beta}}}\end{equation}

This result is valid provided $\alpha > \beta$, which is physically reasonable since the superconducting phase has lower entropy at a given temperature.

\subsubsection*{Part 3: Internal Energies at Finite Temperature}

The internal energy is related to heat capacity by:
\begin{equation}dE = TdS = CdT\end{equation}

For the normal phase, starting from $E_n(0) = E_0$:
\begin{equation}E_n(T) = E_0 + \int_0^T C_n(T')dT' = E_0 + \int_0^T V(\beta T'^3 + \gamma T')dT'\end{equation}

\begin{equation}E_n(T) = E_0 + V\left[\frac{\beta T'^4}{4} + \frac{\gamma T'^2}{2}\right]_0^T\end{equation}

\begin{equation}\boxed{E_n(T) = E_0 + V\left(\frac{\beta T^4}{4} + \frac{\gamma T^2}{2}\right)}\end{equation}

For the superconducting phase, starting from $E_s(0) = E_0 - V\Delta$:
\begin{equation}E_s(T) = E_0 - V\Delta + \int_0^T C_s(T')dT' = E_0 - V\Delta + \int_0^T V\alpha T'^3 dT'\end{equation}

\begin{equation}E_s(T) = E_0 - V\Delta + V\left[\frac{\alpha T'^4}{4}\right]_0^T\end{equation}

\begin{equation}\boxed{E_s(T) = E_0 - V\Delta + \frac{V\alpha T^4}{4}}\end{equation}

\subsubsection*{Part 4: Energy Gap $\Delta$}

At equilibrium, the two phases coexist when their Gibbs free energies are equal. The Gibbs free energy is:
\begin{equation}G = E - TS\end{equation}

At the critical temperature $T_c$ (with zero magnetic field), we have:
\begin{equation}G_n(T_c) = G_s(T_c)\end{equation}

This gives:
\begin{equation}E_n(T_c) - T_c S_n(T_c) = E_s(T_c) - T_c S_s(T_c)\end{equation}

Rearranging:
\begin{equation}E_n(T_c) - E_s(T_c) = T_c[S_n(T_c) - S_s(T_c)]\end{equation}

Substituting our expressions:
\begin{equation}\left[E_0 + V\left(\frac{\beta T_c^4}{4} + \frac{\gamma T_c^2}{2}\right)\right] - \left[E_0 - V\Delta + \frac{V\alpha T_c^4}{4}\right] = T_c\left[V\left(\frac{\beta T_c^3}{3} + \gamma T_c\right) - \frac{V\alpha T_c^3}{3}\right]\end{equation}

Simplifying the left side:
\begin{equation}V\Delta + V\left(\frac{\beta T_c^4}{4} - \frac{\alpha T_c^4}{4} + \frac{\gamma T_c^2}{2}\right) = V\left(\frac{\beta T_c^4}{3} - \frac{\alpha T_c^4}{3} + \gamma T_c^2\right)\end{equation}

Dividing by $V$:
\begin{align}\Delta + \frac{(\beta - \alpha)T_c^4}{4} + \frac{\gamma T_c^2}{2} &= \frac{(\beta - \alpha)T_c^4}{3} + \gamma T_c^2\\
 &= \frac{(\beta - \alpha)T_c^4}{3} - \frac{(\beta - \alpha)T_c^4}{4} + \gamma T_c^2 - \frac{\gamma T_c^2}{2}\\
 &= (\beta - \alpha)T_c^4\left(\frac{1}{3} - \frac{1}{4}\right) + \gamma T_c^2\left(1 - \frac{1}{2}\right)\\
 &= (\beta - \alpha)T_c^4 \cdot \frac{1}{12} + \gamma T_c^2 \cdot \frac{1}{2}\\
&= \frac{(\beta - \alpha)T_c^4}{12} + \frac{\gamma T_c^2}{2}
\end{align}

Using $T_c^2 = \frac{3\gamma}{\alpha - \beta}$:
\begin{align}\Delta &= \frac{(\beta - \alpha)T_c^4}{12} + \frac{\gamma T_c^2}{2} = -\frac{(\alpha - \beta)T_c^4}{12} + \frac{\gamma T_c^2}{2} \\
 &= -\frac{(\alpha - \beta)T_c^2 \cdot T_c^2}{12} + \frac{\gamma T_c^2}{2} = -\frac{3\gamma T_c^2}{12} + \frac{\gamma T_c^2}{2} = -\frac{\gamma T_c^2}{4} + \frac{\gamma T_c^2}{2}
 \end{align}

\begin{equation}\boxed{\Delta = \frac{\gamma T_c^2}{4}}\end{equation}

\subsubsection*{Part 5: Critical Magnetic Field $B_c(T)$}

In the presence of a magnetic field, the appropriate thermodynamic potential is the Gibbs free energy that includes magnetic work. For a system with magnetization $M$, the relevant potential is:
\begin{equation}G = E - TS - BM\end{equation}

For the normal phase with $M_n = 0$:
\begin{equation}G_n(T,B) = E_n(T) - TS_n(T)\end{equation}

For the superconducting phase with $M_s = -\frac{VB}{4\pi}$:
\begin{equation}G_s(T,B) = E_s(T) - TS_s(T) - B \cdot \left(-\frac{VB}{4\pi}\right) = E_s(T) - TS_s(T) + \frac{VB^2}{4\pi}\end{equation}

The phase boundary occurs when $G_n(T,B) = G_s(T,B)$:
\begin{equation}E_n(T) - TS_n(T) = E_s(T) - TS_s(T) + \frac{VB_c^2}{4\pi}\end{equation}

At $B=0$ and $T=T_c$, the phases are in equilibrium, so:
\begin{equation}E_n(T_c) - T_c S_n(T_c) = E_s(T_c) - T_c S_s(T_c)\end{equation}

For general $(T, B_c)$ on the phase boundary:
\begin{equation}[E_n(T) - E_s(T)] - T[S_n(T) - S_s(T)] = \frac{VB_c^2}{4\pi}\end{equation}

Let's define $\Delta G(T) \equiv G_n(T,0) - G_s(T,0)$. At the phase boundary:
\begin{equation}\Delta G(T) = -\frac{VB_c^2}{4\pi}\end{equation}

We need to compute $\Delta G(T)$:
\begin{align}
    \Delta G(T) &= [E_n(T) - E_s(T)] - T[S_n(T) - S_s(T)]\\
    &= \left[V\left(\frac{\beta T^4}{4} + \frac{\gamma T^2}{2}\right) - \left(-V\Delta + \frac{V\alpha T^4}{4}\right)\right] - T\left[V\left(\frac{\beta T^3}{3} + \gamma T\right) - \frac{V\alpha T^3}{3}\right]\\
&= V\Delta + V\left[\frac{(\beta - \alpha)T^4}{4} + \frac{\gamma T^2}{2}\right] - VT\left[\frac{(\beta - \alpha)T^3}{3} + \gamma T\right]\\
&= V\Delta + \frac{V(\beta - \alpha)T^4}{4} + \frac{V\gamma T^2}{2} - \frac{V(\beta - \alpha)T^4}{3} - V\gamma T^2\\
&= V\Delta + V(\beta - \alpha)T^4\left(\frac{1}{4} - \frac{1}{3}\right) + V\gamma T^2\left(\frac{1}{2} - 1\right)\\
&= V\Delta - \frac{V(\beta - \alpha)T^4}{12} - \frac{V\gamma T^2}{2}\end{align}

Using $\Delta = \frac{\gamma T_c^2}{4}$ and $(\alpha - \beta) = \frac{3\gamma}{T_c^2}$:
\begin{align}
    \Delta G(T) &= \frac{V\gamma T_c^2}{4} - \frac{V \cdot 3\gamma \cdot T^4}{12T_c^2} - \frac{V\gamma T^2}{2}\\
&= \frac{V\gamma T_c^2}{4} - \frac{V\gamma T^4}{4T_c^2} - \frac{V\gamma T^2}{2}\\
&= \frac{V\gamma}{4}\left(T_c^2 - \frac{T^4}{T_c^2} - 2T^2\right)\\
&= \frac{V\gamma}{4T_c^2}\left(T_c^4 - T^4 - 2T^2T_c^2\right) = \frac{V\gamma}{4T_c^2}\left(T_c^2 - T^2\right)^2
\end{align}

Setting this equal to $-\frac{VB_c^2}{4\pi}$:
\begin{equation}\frac{V\gamma(T_c^2 - T^2)^2}{4T_c^2} = -\frac{VB_c^2}{4\pi}\end{equation}

Wait, this gives a negative $B_c^2$, which is unphysical. Let me reconsider the sign. The superconducting phase should have lower free energy at $B=0$, so $\Delta G(T) < 0$ for $T < T_c$. Actually:

\begin{align}
    -\Delta G(T) &= \frac{VB_c^2}{4\pi}\\
\frac{VB_c^2}{4\pi} &= \frac{V\gamma(T_c^2 - T^2)^2}{4T_c^2}\\
B_c^2 &= \frac{\pi\gamma(T_c^2 - T^2)^2}{T_c^2} = \frac{\pi\gamma T_c^2(T_c^2 - T^2)^2}{T_c^4} = \frac{\pi\gamma T_c^2}{T_c^4}\left(1 - \frac{T^2}{T_c^2}\right)^2T_c^4
\end{align}

Actually, let me factor more carefully:
\begin{equation}B_c^2 = \pi\gamma T_c^2\left(1 - \frac{T^2}{T_c^2}\right)^2\end{equation}

Therefore:
\begin{equation}\boxed{B_c(T) = B_0\left(1 - \frac{T^2}{T_c^2}\right)}\end{equation}

where:
\begin{equation}\boxed{B_0 = \sqrt{\pi\gamma T_c^2} = \sqrt{\pi\gamma} \cdot T_c}\end{equation}

This parabolic dependence of the critical field on temperature is characteristic of type-I superconductors and is well-confirmed experimentally.
\section*{Problem 4}
\subsection*{Problem Statement}
\textit{Photon gas Carnot cycle:} the aim of this problem is to obtain the black-body radiation relation, $E(T,V)\propto VT^4$, starting from the equation of state, by performing an infinitesimal Carnot cycle on the photon gas.
\begin{itemize}
    \item Express the work done, $W$, in the above cycle, in terms of $dV$ and $dP$.
    \item Express the heat absorbed $Q$, in expanding the gas along an isotherm, in terms of $P$, $dV$ and an appropriate derivative of $E(T,V)$.
    \item Using the efficiency of the Carnot cycle, relate the above expressions for $W$, and $Q$ to $T$ and $dT$.
    \item Observations indicate that the pressure of the photon gas is given by $P=AT^4$, where $A = \pi^2 k_B^2 / 45(\hbar c)^3$ is a constant. Use this information to obtain $E(T,V)$, assuming  $E(0, V)=0$.
    \item Find the relation describing adiabatic paths in the above cycle.
\end{itemize}

\subsection*{Solution}

This problem uses the powerful method of analyzing an infinitesimal Carnot cycle to derive thermodynamic relations for a photon gas (blackbody radiation). The Carnot cycle consists of two isothermal processes and two adiabatic processes, and its efficiency provides a fundamental constraint that connects various thermodynamic quantities.

\subsubsection*{Part 1: Work Done in the Infinitesimal Carnot Cycle}

Consider an infinitesimal Carnot cycle operating between temperatures $T$ and $T + dT$, with volume changes between $V$ and $V + dV$. The cycle consists of four processes:

\begin{enumerate}
    \item Isothermal expansion at temperature $T$ from $(V, P)$ to $(V + dV, P')$
    \item Adiabatic expansion from $(V + dV, T)$ to $(V + dV, T + dT)$ with pressure change
    \item Isothermal compression at temperature $T + dT$ from $(V + dV, P'')$ to $(V, P + dP)$
    \item Adiabatic compression from $(V, T + dT)$ back to $(V, T)$
\end{enumerate}

The work done in an infinitesimal cycle can be computed as the area enclosed by the cycle in the $P$-$V$ diagram. For an infinitesimal rectangular cycle, this area is approximately:
\begin{equation}W = \oint P \, dV\end{equation}

The cycle forms a small parallelogram in the $P$-$V$ plane. The net work is the difference between the work done during expansion and compression processes. For an infinitesimal cycle, this gives:
\begin{equation}\boxed{W = dP \cdot dV}\end{equation}

This represents the area of the infinitesimal parallelogram enclosed by the cycle, where $dP$ is the pressure difference between the two isotherms and $dV$ is the volume difference.

\subsubsection*{Part 2: Heat Absorbed During Isothermal Expansion}

During the isothermal expansion at temperature $T$ from volume $V$ to $V + dV$, the first law of thermodynamics states:
\begin{equation}dE = \delta Q - \delta W\end{equation}

For an isothermal process at constant temperature $T$, the change in internal energy is:
\begin{equation}dE\big|_T = \left(\frac{\partial E}{\partial V}\right)_T dV\end{equation}

The work done during this isothermal expansion is:
\begin{equation}\delta W = P \, dV\end{equation}

Therefore, the heat absorbed during the isothermal expansion is:
\begin{equation}Q = dE\big|_T + P \, dV = \left(\frac{\partial E}{\partial V}\right)_T dV + P \, dV\end{equation}

\begin{equation}\boxed{Q = \left[P + \left(\frac{\partial E}{\partial V}\right)_T\right] dV}\end{equation}

This expression shows that the heat absorbed depends not only on the $PdV$ work but also on how the internal energy changes with volume at constant temperature.

\subsubsection*{Part 3: Relating Work, Heat, and Efficiency}

The efficiency of a Carnot cycle operating between temperatures $T$ and $T + dT$ is:
\begin{equation}\eta = \frac{W}{Q_h} = \frac{T + dT - T}{T + dT} = \frac{dT}{T + dT} \approx \frac{dT}{T}\end{equation}

where $Q_h = Q$ is the heat absorbed at the higher temperature during isothermal expansion. The work done is:
\begin{equation}W = \eta \cdot Q = \frac{dT}{T} \cdot Q\end{equation}

Substituting our expressions for $W$ and $Q$:
\begin{equation}dP \cdot dV = \frac{dT}{T} \cdot \left[P + \left(\frac{\partial E}{\partial V}\right)_T\right] dV\end{equation}

Dividing both sides by $dV$ (assuming $dV \neq 0$):
\begin{equation}dP = \frac{dT}{T} \cdot \left[P + \left(\frac{\partial E}{\partial V}\right)_T\right]\end{equation}

Since we're considering changes along the cycle, $dP$ represents the change in pressure between the two isotherms at the same volume, so:
\begin{equation}\left(\frac{\partial P}{\partial T}\right)_V dT = \frac{dT}{T} \cdot \left[P + \left(\frac{\partial E}{\partial V}\right)_T\right]\end{equation}

Dividing by $dT$:
\begin{equation}\boxed{\left(\frac{\partial P}{\partial T}\right)_V = \frac{1}{T}\left[P + \left(\frac{\partial E}{\partial V}\right)_T\right]}\end{equation}

This is actually a fundamental thermodynamic relation that can also be derived from Maxwell relations. Rearranging:
\begin{equation}T\left(\frac{\partial P}{\partial T}\right)_V = P + \left(\frac{\partial E}{\partial V}\right)_T\end{equation}

\subsubsection*{Part 4: Deriving the Energy Relation $E(T,V)$}

Given that the pressure of the photon gas is:
\begin{equation}P = AT^4\end{equation}

We can compute:
\begin{equation}\left(\frac{\partial P}{\partial T}\right)_V = 4AT^3\end{equation}

Substituting into our relation from Part 3:
\begin{align}
    T \cdot 4AT^3 &= AT^4 + \left(\frac{\partial E}{\partial V}\right)_T\\
4AT^4 &= AT^4 + \left(\frac{\partial E}{\partial V}\right)_T\\
\left(\frac{\partial E}{\partial V}\right)_T &= 3AT^4\end{align}

This tells us that at constant temperature, the internal energy is proportional to volume:
\begin{equation}E(T,V) = 3AT^4 \cdot V + f(T)\end{equation}

where $f(T)$ is an arbitrary function of temperature only. However, for a photon gas, the energy must be extensive in volume, meaning that when $V = 0$, we must have $E = 0$ regardless of temperature (there are no photons in zero volume). This requires $f(T) = 0$.

Using the boundary condition $E(0, V) = 0$, we confirm that there are no additional temperature-dependent constants. Therefore:
\begin{equation}\boxed{E(T,V) = 3AT^4 V}\end{equation}

Since $A = \frac{\pi^2 k_B^4}{45(\hbar c)^3}$, we can write:
\begin{equation}E(T,V) = \frac{3\pi^2 k_B^4}{45(\hbar c)^3} VT^4 = \frac{\pi^2 k_B^4}{15(\hbar c)^3} VT^4\end{equation}

This is the Stefan-Boltzmann law for blackbody radiation, showing that the energy density is proportional to $T^4$.

\subsubsection*{Part 5: Adiabatic Path Relation}

For an adiabatic process, $\delta Q = 0$, so by the first law:
\begin{equation}dE = -P \, dV\end{equation}

The total differential of $E(T,V) = 3AT^4V$ is:
\begin{equation}dE = \left(\frac{\partial E}{\partial T}\right)_V dT + \left(\frac{\partial E}{\partial V}\right)_T dV = 12AT^3V \, dT + 3AT^4 \, dV\end{equation}

Setting $dE = -P \, dV = -AT^4 \, dV$:
\begin{align}
    12AT^3V \, dT + 3AT^4 \, dV &= -AT^4 \, dV\\
12AT^3V \, dT &= -AT^4 \, dV - 3AT^4 \, dV = -4AT^4 \, dV\end{align}

Dividing by $AT^3$:
\begin{equation}12V \, dT = -4T \, dV\end{equation}

\begin{equation}3V \, dT = -T \, dV\end{equation}

Separating variables:
\begin{equation}\frac{dT}{T} = -\frac{dV}{3V}\end{equation}

Integrating both sides:
\begin{align}
    \ln T &= -\frac{1}{3}\ln V + \text{constant}\\
\ln T + \frac{1}{3}\ln V &= \text{constant}\\
\ln(T V^{1/3}) &= \text{constant}
\end{align}

Therefore, the adiabatic relation is:
\begin{equation}\boxed{TV^{1/3} = \text{constant}}\end{equation}

or equivalently:
\begin{equation}\boxed{T^3V = \text{constant}}\end{equation}

This is the adiabatic equation for a photon gas, showing that during an adiabatic expansion (increase in $V$), the temperature decreases as $T \propto V^{-1/3}$. This relation is fundamental in cosmology, describing how the temperature of the cosmic microwave background radiation decreases as the universe expands.
\section*{Problem 5}
\subsection*{Problem Statement}
\textit{Irreversible Processes}
\begin{itemize}
    \item Consider two substances, initially at temperatures $T^0_1$ and $T^0_2$, coming to equilibrium at a final temperature $T_f$ through heat exchange. By relating the direction of heat flow to the temperature difference, show that the change in the total entropy which can be written as:
    \begin{equation}
        \Delta S = \Delta S_1 + \Delta S_2 \geq \int_{T_1^0}^{T_f} \frac{\delta Q_1}{T_1} + \int_{T_2^0}^{T_f}\frac{\delta Q_2}{T_2} = \int \frac{T_1 - T_2}{T_1T_2}\delta Q
    \end{equation}
    must be positive. This is an example of the more general condition that ``in a closed system, equilibrium is characterized by the maximum value of entropy $S$''
    \item Now consider a gas with adjustable volume $V$, and diathermal walls, embedded in a heat bath of constant temperature $T$, and fixed pressure $P$. The change in the entropy of the bath is given by:
    \begin{equation}
        \Delta S_{\text{bath}} = \frac{\Delta Q_{\text{bath}}}{T} = -\frac{\Delta Q_{\text{gas}}}{T} = \frac{1}{T}\left(\Delta E_{\text{gas}} + P \Delta V_{\text{gas}}\right).
    \end{equation}
    \item By considering the change in entropy of the combined system establish that ``the equilibrium of a gas at fixed $T$ and $P$ is characterized by the minimum of the Gibbs free energy $G= E + PV - TS$.''
\end{itemize}

\subsection*{Solution}

This problem explores fundamental principles of thermodynamics related to irreversible processes and the conditions for equilibrium. We will demonstrate how entropy maximization in isolated systems leads to the minimization of appropriate free energies under different constraints.

\subsubsection*{Part 1: Entropy Change in Heat Exchange Between Two Substances}

Consider two substances initially at temperatures $T_1^0$ and $T_2^0$ that exchange heat until they reach thermal equilibrium at a common final temperature $T_f$. Without loss of generality, assume $T_1^0 > T_2^0$, so heat flows from substance 1 to substance 2.

For each substance, the entropy change can be written as:
\begin{equation}\Delta S_1 = \int_{T_1^0}^{T_f} \frac{\delta Q_1}{T_1}, \quad \Delta S_2 = \int_{T_2^0}^{T_f} \frac{\delta Q_2}{T_2}\end{equation}

where $\delta Q_1$ and $\delta Q_2$ are the infinitesimal heat transfers to substances 1 and 2, respectively. Since the two substances form an isolated system, energy conservation requires:
\begin{equation}\delta Q_1 + \delta Q_2 = 0 \quad \Rightarrow \quad \delta Q_2 = -\delta Q_1\end{equation}

Let us denote $\delta Q \equiv \delta Q_2 = -\delta Q_1$ as the heat flowing from substance 1 to substance 2. Since heat flows from hot to cold, we have $\delta Q > 0$ when $T_1 > T_2$.

The total entropy change of the combined system is:
\begin{equation}\Delta S = \Delta S_1 + \Delta S_2 = \int_{T_1^0}^{T_f} \frac{\delta Q_1}{T_1} + \int_{T_2^0}^{T_f} \frac{\delta Q_2}{T_2}\end{equation}

Substituting $\delta Q_1 = -\delta Q$ and $\delta Q_2 = \delta Q$:
\begin{equation}\Delta S = \int_{T_1^0}^{T_f} \frac{-\delta Q}{T_1} + \int_{T_2^0}^{T_f} \frac{\delta Q}{T_2} = \int \left(\frac{1}{T_2} - \frac{1}{T_1}\right) \delta Q\end{equation}

Rewriting with a common denominator:
\begin{equation}\Delta S = \int \frac{T_1 - T_2}{T_1 T_2} \delta Q\end{equation}

Now we demonstrate that $\Delta S > 0$ for an irreversible process. The key physical insight is that heat flows spontaneously from hot to cold, meaning:
\begin{itemize}
\item When $T_1 > T_2$, heat flows from 1 to 2, so $\delta Q > 0$
\item When $T_1 < T_2$, heat flows from 2 to 1, so $\delta Q < 0$
\end{itemize}
In both cases, the product $(T_1 - T_2)\delta Q$ has the same sign:
\begin{itemize}
\item If $T_1 > T_2$: both factors positive, so $(T_1 - T_2)\delta Q > 0$
\item If $T_1 < T_2$: both factors negative, so $(T_1 - T_2)\delta Q > 0$
\end{itemize}
Since $T_1 T_2 > 0$ always, we have:
\begin{equation}\frac{T_1 - T_2}{T_1 T_2} \delta Q > 0\end{equation}

Therefore:
\begin{equation}\boxed{\Delta S = \int \frac{T_1 - T_2}{T_1 T_2} \delta Q > 0}\end{equation}

The entropy change is strictly positive for any irreversible heat transfer process. Only in the reversible limit, where the temperature difference is infinitesimal at each step, does $\Delta S \to 0$. This demonstrates that in an isolated system undergoing spontaneous processes, the entropy always increases until it reaches a maximum at equilibrium. The condition for equilibrium is thus:
\begin{equation}\delta S = 0 \quad \text{(maximum entropy)}\end{equation}

This is the principle of entropy maximization: **in a closed system with fixed energy and volume, equilibrium is characterized by the maximum value of entropy**.

\subsubsection*{Part 2: Entropy Change of the Heat Bath}

Consider a gas with adjustable volume $V$ enclosed by diathermal (heat-conducting) walls, in contact with a heat bath at constant temperature $T$ and subject to constant external pressure $P$. The gas can exchange heat with the bath and perform work by changing its volume.

The heat bath absorbs heat $\Delta Q_{\text{bath}}$ from the gas (or vice versa). Since the bath remains at constant temperature $T$, its entropy change is:
\begin{equation}\Delta S_{\text{bath}} = \frac{\Delta Q_{\text{bath}}}{T}\end{equation}

By energy conservation, the heat absorbed by the bath equals the heat released by the gas:
\begin{equation}\Delta Q_{\text{bath}} = -\Delta Q_{\text{gas}}\end{equation}

From the first law of thermodynamics applied to the gas:
\begin{equation}\Delta E_{\text{gas}} = \Delta Q_{\text{gas}} - W_{\text{by gas}}\end{equation}

where the work done by the gas against the constant external pressure is:
\begin{equation}W_{\text{by gas}} = P \Delta V_{\text{gas}}\end{equation}

Therefore:
\begin{equation}\Delta Q_{\text{gas}} = \Delta E_{\text{gas}} + P \Delta V_{\text{gas}}\end{equation}

Substituting into the expression for the bath's entropy change:
\begin{equation}\boxed{\Delta S_{\text{bath}} = -\frac{\Delta Q_{\text{gas}}}{T} = -\frac{1}{T}\left(\Delta E_{\text{gas}} + P \Delta V_{\text{gas}}\right)}\end{equation}

This result shows that the entropy change of the bath is determined by the internal energy change and the $PV$ work performed by the gas.

\subsubsection*{Part 3: Gibbs Free Energy Minimization Principle}

The total entropy change of the combined system (gas plus bath) must be non-negative by the second law:
\begin{equation}\Delta S_{\text{total}} = \Delta S_{\text{gas}} + \Delta S_{\text{bath}} \geq 0\end{equation}

Substituting our expression for $\Delta S_{\text{bath}}$:
\begin{equation}\Delta S_{\text{gas}} - \frac{1}{T}\left(\Delta E_{\text{gas}} + P \Delta V_{\text{gas}}\right) \geq 0\end{equation}

Multiplying through by $-T$ (which reverses the inequality):
\begin{equation}-T\Delta S_{\text{gas}} + \Delta E_{\text{gas}} + P \Delta V_{\text{gas}} \leq 0\end{equation}

Rearranging:
\begin{equation}\Delta E_{\text{gas}} + P \Delta V_{\text{gas}} - T\Delta S_{\text{gas}} \leq 0\end{equation}

We can rewrite this in terms of the change in a thermodynamic potential. Define the **Gibbs free energy** as:
\begin{equation}G \equiv E + PV - TS\end{equation}

For a process at constant temperature $T$ and constant pressure $P$:
\begin{equation}\Delta G = \Delta E + P\Delta V + V\Delta P - T\Delta S - S\Delta T\end{equation}

Since $\Delta P = 0$ and $\Delta T = 0$:
\begin{equation}\Delta G = \Delta E + P\Delta V - T\Delta S\end{equation}

Therefore, our inequality becomes:
\begin{equation}\boxed{\Delta G_{\text{gas}} \leq 0}\end{equation}

This fundamental result states that **for any spontaneous process occurring in a system at constant temperature and pressure, the Gibbs free energy must decrease or remain constant**. At equilibrium, where no further spontaneous changes occur:
\begin{equation}\Delta G = 0 \quad \text{(at equilibrium)}\end{equation}

or equivalently, $G$ is at a minimum with respect to all possible variations:
\begin{equation}\delta G = 0 \quad \text{and} \quad \delta^2 G > 0\end{equation}

This establishes that **the equilibrium of a gas at fixed $T$ and $P$ is characterized by the minimum of the Gibbs free energy** $G = E + PV - TS$.

\subsubsection*{Physical Interpretation}

These results illustrate a profound principle of thermodynamics: the appropriate equilibrium criterion depends on the constraints imposed on the system:
\begin{itemize}
    \item \textbf{Isolated system} (constant $E$, $V$, $N$): Equilibrium at \textbf{maximum $S$}
    \item \textbf{Constant $T$, $V$, $N$}: Equilibrium at \textbf{minimum $F = E - TS$} (Helmholtz free energy)
\item \textbf{Constant $T$, $P$, $N$}: Equilibrium at \textbf{minimum $G = E + PV - TS$} (Gibbs free energy)
\item \textbf{Constant $S$, $P$, $N$}: Equilibrium at \textbf{minimum $H = E + PV$} (Enthalpy)
\end{itemize}
Each free energy represents the appropriate thermodynamic potential whose minimization determines equilibrium under the corresponding constraints. The Gibbs free energy is particularly important in chemistry and materials science because most experiments are conducted at constant temperature and pressure (atmospheric conditions).
\end{document}
